\chapter{Introduction}
\label{ch:introduction}

% A thesis introduction normally states the context, the precise question, the
% contribution, and a short reading guide. For a long report, four to six
% pages are often enough; this is planning advice, not an LMFI rule.

Directed graphs model one-way relations: a dependency may point from one task
to another, or a transition may lead from one program state to its successor.
A first question is whether a vertex can be reached from another by following
zero or more directed edges. This elementary problem is useful here because it
connects mathematical definitions, proofs, algorithms, figures, tables, and
references in one coherent example. Standard introductions can be found in
graph theory and algorithms texts such as \cite{diestel2017,cormen2022}.

\section{Question and contribution}

The example answers the following question: how can reachability in a finite
directed graph be defined, proved to be transitive, and computed from an
adjacency matrix? The contribution is deliberately modest. We formalise the
relation in \cref{ch:background}, compute it in \cref{ch:main-work}, and inspect
a concrete instance in \cref{ch:evaluation}.

\section{Structure of the example}

The separation between background, main work, and evaluation demonstrates a
maintainable chapter structure. A theoretical thesis may split the main work
into several result chapters; an applied thesis may instead split evaluation
into protocol, results, and discussion. The source files are designed to make
either change local and explicit.
